\documentclass[11pt]{article}
\usepackage[T1]{fontenc}
\usepackage[utf8]{inputenc}
\usepackage{lmodern,amsmath,amssymb,amsthm,mathtools,booktabs,array}
\usepackage[margin=1in]{geometry}
\usepackage[colorlinks=true,linkcolor=blue,citecolor=blue,urlcolor=blue]{hyperref}
\usepackage{microtype}
\newtheorem{theorem}{Theorem}[section]
\newtheorem{lemma}[theorem]{Lemma}
\newtheorem{proposition}[theorem]{Proposition}
\newtheorem{corollary}[theorem]{Corollary}
\theoremstyle{definition}
\newtheorem{definition}[theorem]{Definition}
\theoremstyle{remark}
\newtheorem{remark}[theorem]{Remark}
\newcommand{\R}{\mathbb R}
\newcommand{\N}{\mathbb N}
\newcommand{\one}{\mathbf 1}
\DeclareMathOperator*{\argmin}{arg\,min}
\title{The Funding Cost of Preserving Baseline Allocations\\under Resource Withholding}
\author{Working paper candidate}
\date{24 October 2026}
\begin{document}
\maketitle
\begin{abstract}
A coalition-covering allocation at each resource report need not discourage an
owner from withholding resources. We study repairs that make each owner's payment
nondecreasing in its own reported bundle while preserving every baseline payment
at every profile. The coordinatewise least such repair is the own-report monotone
envelope. For finite min-linear production technologies with fixed owners and a
continuous report domain, every efficient-core baseline selection can require
funding linear in the number of owners after repair. The same conclusion holds
for every baseline with a fixed funding factor below two. Our principal result
extends this linear lower order to every fixed finite funding cap and every fixed
strictly convex power objective used to select the baseline. The elementary
linear upper bound is therefore sharp in owner order along an unbounded sequence.
Without a funding cap, the quadratic baseline instead admits envelope lower
bounds exceeding every constant multiple of the square root of the owner count.
The constructions keep one technology and the same actual owners throughout all
report comparisons. These are limits on baseline-preserving incentive repair;
such growing lower bounds are not established here for arbitrary truthful
payment mechanisms.
\end{abstract}

\section{Introduction}
Resource pooling creates two distinct allocation problems. At a given contribution
profile, a payment vector should cover what every coalition can produce on its
own. Across profiles, an owner should not receive a larger payment by reporting
a smaller bundle. The first requirement concerns the core or an externally funded
extension of it. The second concerns the incentives induced by a rule that chooses
one allocation at every possible report. Solving the first problem independently
at each profile need not solve the second.

We examine a specific and explicit restriction on repairs. A baseline rule has
already assigned a payment to every owner at every profile. A revised rule may
add payments, but may not reduce any of these baseline allocations. This is a
conditional design principle: it can represent a commitment to an existing
allocation schedule. We do not assert that an institution has adopted the rules
studied here, or that this commitment is required by the unrestricted mechanism
problem. The question is what funding is necessary if that commitment is imposed.
It is stronger than preserving the allocation at one observed profile: the floor
applies to the entire report-contingent schedule.

For a baseline $q$, the least permissible incentive repair assigns owner $i$
the payment
\[
 E_i^q(b)=\sup_{0\leq a_i\leq b_i}q_i(a_i,b_{-i}).
\]
This elementary order construction preserves coalition coverage and is monotone
in the owner's own bundle. Every other own-monotone rule dominating $q$ at all
profiles must dominate $E^q$ as well. Thus a lower bound for the envelope is a
lower bound for every repair satisfying the same baseline-floor contract. The
operator itself is not offered as a new mechanism-design technique.

Our first result is an elementary obstruction that does not depend on a convex
objective. For every selection from the efficient core, some profile with $N$
actual owners requires envelope funding at least $(N+6)/8$ times production,
along the even owner counts in the construction. An extension covers arbitrary
baseline rules whose total payments are bounded by $B$ times production for any
fixed $1\leq B<2$. These results quantify an incompatibility between tight
baseline budgets and incentives even when the selection is discontinuous,
asymmetric, or otherwise unrestricted.

The principal construction concerns a broader family of proposed baseline rules:
minimize $\sum_i x_i^r$ over nonnegative coalition covers, subject to a total-payment
cap $C$ times production. For every fixed finite $C\geq1$ and fixed $r>1$, the
envelope needs $\Omega_C(N)$ funding along an unbounded sequence. At $C=2$ the
explicit coefficient is $1/11$, up to one owner. The upper bound $NC$ is immediate,
so the owner order is sharp in this sequence sense. We do not obtain matching
constants, a lower bound at every prescribed owner count, or a bound with the
resource dimension held fixed. The parameters $C$ and $r$ are fixed before the
sequence is taken; its onset may depend on $r$.

The uncapped quadratic rule requires a separate analysis. Its envelope exceeds
every universal constant times $\sqrt N$; the explicit sequence gives a factor
of order at least $\sqrt N(\log N/\log\log N)^{1/3}$. The fixed-cap theorem does
not imply this statement: a cap large enough to be inactive can grow with $N$.
The two results belong to one account of the funding cost of preserving convexly
chosen allocations. A restricted common-veto refinement is retained with its
precise class assumptions, rather than treated as another paper.

The technical work is to force many owners' distinct underreports to create large
baseline payments while preserving one technology and one owner set. An exact
row-mixture construction supplies a convex optimality certificate; a control-resource
realization implements the required row changes as actual unilateral withholding.
The appendix gives the transport argument for every required cut. Finite arithmetic
checks accompany the proof but do not enumerate the enormous final instances or
replace the all-parameter existence argument.

\subsection{Relationship to existing work and scope of contribution}
Minimum-norm selection from the efficient core is classical. Arin, Kuipers and
Vermeulen \cite{AKV2008}, Definition~1, define the least-squares core allocation
by precisely this criterion. At $C=1,r=2$ our baseline is their selection. Under
efficiency, minimizing squared norm is equivalent to minimizing squared distance
to equal division. That equivalence need not hold when a subsidy makes total
payments variable. We treat the capped and uncapped power variants as specified
design choices, without claiming their institutional adoption.

Static minimum-norm coalition covers and the associated projection inequalities
also appear in the critical-threshold analysis of Hof et al.\ \cite{Hof2018}.
The blocking-polyhedron optimality identities used here have close classical
antecedents, including Abdi, Cornu\'{e}jols and Lee \cite{ACL2019}. Our candidate
contribution is the quantitative aggregation of payments across genuine
unilateral report changes, not the static projection identity or its multiplier
representation.

Representation must likewise be separated from incentives. Kalai and Zemel
\cite{KZ1980} represent totally balanced games through nonnegative additive forms
and flow games. Our construction uses an elementary control-coordinate device
to realize several linked report profiles simultaneously with fixed owners.
We give that device explicitly and do not claim a new general representation
theorem. The minimum-norm constructions of Kultti and Salonen \cite{KS2005}
use different objects and objectives, including approximation of a characteristic
function by additive games; their terminology alone does not identify our baseline.

Core selection and monotonicity have a substantial literature. Arin
\cite{Arin2011} uses precisely the incident-coalition direction induced by an
owner increasing its bundle: values of coalitions containing that owner increase,
and other coalition values stay fixed. His impossibility theorem and linked-game
compatibility conditions are close qualitative antecedents. They do not state
an owner-count funding rate or the fixed-cap, fixed-power theorem proved here.
Hougaard, Peleg and \O sterdal \cite{HPO2003} study strictly concave objectives
over efficient cores and prove coalitional monotonicity on convex games. Our
pool examples are not convex games; a direct supermodularity violation is given
below. We claim no new efficient-core objective or general qualitative
incompatibility. Izquierdo and
Rafels \cite{IR2017} study resource incentives in scalar co-investment problems
with increasing average returns. Lahiri \cite{Lahiri2005} gives market-equilibrium
manipulation examples in production economies, including a constant-returns
example. Resource withholding itself is therefore also an established concern;
our question is the quantitative funding cost of a specified payment-floor repair.
Dietzenbacher and Do\u{g}an \cite{DD2025}
characterize the compatibility of egalitarian core selection and population
monotonicity. Their changing population and subgame comparisons differ from our
fixed-owner bundle reports, but the incompatibility examples are a relevant
structural precedent. Our elementary efficient-core obstruction should be read
as a quantitative resource-withholding separation, not as discovery of a general
conflict between stability and monotonicity.

The cost-of-stability framework of Bachrach et al.\ \cite{CostStability2009}
asks how much external money is needed to stabilize a given game. Every game in
our production model already has an efficient core allocation. Our extra funding
arises from preserving allocations across reports while imposing incentives.
Likewise, monotonicity of a quadratic regularization path on a fixed polytope
\cite{GNR2025} concerns a different varying parameter from the incident coalition
constraints changed by one owner's report.

Accordingly, the proposed contribution is the explicit owner-count lower bounds
for this additional baseline-floor contract, especially the fixed-cap,
fixed-power theorem and the uncapped quadratic separation. We make no claim to
resolve the unrestricted funding problem, establish coalition-strategyproofness,
handle private opportunity costs, or provide a computationally efficient mechanism.
Historical priority of these precise quantitative statements remains a question
for specialist assessment; the source comparison accompanying this draft states
which original versions and theorem passages were checked.

